\section{Introduction}

Prior work established that reviewer profile depth is the strongest predictor of evaluation quality in AI expert panel assessments of puzzle hunt puzzles~\cite{paper1}. A bare reviewer persona — ``you are a puzzle designer evaluating this puzzle'' — produces serviceable but generic critiques; a fully elaborated profile incorporating domain knowledge, design philosophy, and a named lens vocabulary produces critiques that identify structural failures the bare condition misses entirely. The gradient is non-monotonic: domain expertise without a design framework adds negligible quality over bare baseline, while a philosophical worldview alone captures most of the value of the full profile. These findings motivate a natural question with a non-obvious answer: does the same depth gradient hold when the profile guides \emph{creation} rather than critique?

The question is non-trivial because reviewing and authoring are cognitively distinct activities. A reviewer applies a lens to finished material; the lens is terminal, in the sense that it requires completed artifacts to have anything to evaluate. An author must generate material that does not yet exist; the profile must be generative, producing constraints that can be satisfied during construction rather than checks that can only be applied after the fact. A reviewer profile that asks ``does the visual hierarchy resolve cleanly?'' presupposes a layout to inspect; the same question offered as an authoring instruction is incoherent when the format is a text acrostic. The transfer from reviewing expertise to authoring expertise is therefore not guaranteed — and, we show, depends critically on the \emph{type} of lens the reviewer applies.

We conducted a 7-condition authoring ablation (A0 through A6) that mirrors the reviewer ablation of Paper~I exactly: A0 provides only a bare authoring persona, A1 adds task specification, A2 adds domain world data, A3 adds design principles, A4 adds an authoring philosophy, A5 adds a full authoring profile, and A6 repurposes a C8-tier reviewing profile as an authoring guide. All seven conditions produced puzzles evaluated by the same three-reviewer C8 panel (Katz, Snyder, Young) on a weighted 45-point rubric. We then extended the study with a 3-designer $\times$ 2-mode comparison: Snyder, Katz, and Young each authored puzzles in both their full authoring profile (A5) and their own reviewing profile repurposed for authoring (A6), yielding six additional evaluated puzzles. The evaluation session was run with all 13 puzzles read before scoring to eliminate calibration drift.

Four findings emerged. First, the authoring quality gradient mirrors the reviewing quality gradient in shape: improvement is real but non-monotonic, with the largest single jump occurring at the principle-introduction boundary (A2 to A3, +2.7 points) and a plateau-and-retreat at the top of the Snyder ablation group (A5 at 32.0/45, A6 at 31.3/45). This symmetric gradient confirms that the mechanisms governing profile depth are not domain-specific to evaluation: the same structural properties of a profile — framework versus content, worldview versus checklist — determine quality in both directions. Second, A2 demonstrates an inversion that has practical implications: domain expertise without a design framework (28.3/45) is nearly indistinguishable from bare task specification (27.7/45). A data-rich profile that lacks generative structure produces technically accurate but experientially flat output — the clues read as database rows rather than puzzle constructions. Third, authoring worldview alone (A4, 30.7/45) is within 1.3 points of the full authoring profile (A5, 32.0/45), a gap smaller than any other adjacent condition pair except A0–A1. The full profile's marginal contribution over worldview is verification rigor — elimination of competing candidates, uniqueness checking — rather than the elevation of mechanism quality, confirming that framework matters more than completeness. Fourth, and most distinctively, the A5-to-A6 transfer result depends entirely on lens type: Katz's operational lens (named testable operations that invert cleanly to authoring constraints) improved by +1.6 points in A6 mode; Young's perceptual lens (checks requiring finished visual or experiential material) improved by +3.7 points — not because perceptual checks transferred, but because Young's reviewer-as-author exercise forced a novel solution within the format's constraints, replacing visual grammar with second-person narrative grammar; and Snyder's construction lens changed by only $-0.7$ points because his reviewing and authoring profiles are already functionally equivalent.

The lens-type result is the paper's primary theoretical contribution. We propose a three-category transfer model: \emph{construction lenses} (Snyder) are already generative — the reviewing and authoring profiles collapse into one another; \emph{operational lenses} (Katz) invert cleanly — each named test has a corresponding authoring-time constraint that can be applied before the puzzle exists; \emph{perceptual lenses} (Young) are format-dependent — they succeed when the output format can supply the required material, and produce either failure or forced workarounds when it cannot. The model predicts which reviewer profiles can be safely repurposed for authoring tasks without redesign, which require inversion, and which require replacement.

The testbed is the same used in Paper~I: puzzle hunt puzzles in the Age of Empires~II domain, with fixed answer word (\textsc{tower}), fixed mechanism (five-clue first-letter acrostic), and locked world data. The controlled format enables clean ablation comparisons while exposing a format ceiling: all A-series puzzles cap at Riven Standard $\leq 3$ because the first-letter acrostic mechanism is imported from puzzle hunt tradition rather than native to the AoE2 domain. This ceiling is itself informative — it establishes that profile depth raises the floor on execution quality within a fixed mechanism, while mechanism choice determines the ceiling, a distinction with direct implications for how authoring profiles should be deployed in practice.

The remainder of this paper is organized as follows. Section~2 situates the work relative to prior research on persona prompting, reviewer profile effects, and the authoring/reviewing distinction in human creativity research. Section~3 describes the authoring ablation design, the designer comparison methodology, and the evaluation rubric. Section~4 reports quantitative results across all 13 evaluated puzzles and introduces the six-framework taxonomy surfaced by richer profiles. Section~5 interprets the symmetric gradient, the lens-type transfer model, and the format ceiling. Section~6 concludes with practical recommendations for constructing authoring profiles from reviewing profiles, and identifies mechanism novelty as the open problem motivating future work.
